\section{Conventions and preliminary facts}\label{sec:conventions}

We fix the entropy normalizations and matrix convergence facts used at
both noise scales. The matrix inputs may be random tuples supported on
bounded microstate sets.

All noncommutative variables below belong to tracial von Neumann algebras
with faithful normal tracial state $\tau$. Write $\|a\|_2^2=\tau(a^*a)$.
A standard semicircular family has variance one in each coordinate and
is free. Every noise family is free from the tuple it perturbs.

For a bounded self-adjoint $n$-tuple $X$, let
$\Gamma_R(X;m,d,\varepsilon)$ consist of self-adjoint $d\times d$ matrix
tuples of coordinate operator norms at most $R$ whose normalized traces
of all words of length at most $m$ differ from those of $X$ by less than
$\varepsilon$. With volume taken in the real Euclidean space defined by
unnormalized Hilbert--Schmidt inner product, set
\[
 \chi_R(X)=\inf_{m,\varepsilon}\limsup_{d\to\infty}
 \left[d^{-2}\log\Vol\Gamma_R(X;m,d,\varepsilon)
       +\frac n2\log d\right],\qquad
 \chi(X)=\sup_R\chi_R(X).
\]
For a tuple $Z$, its conjugate variables $\xi_i\in L^2(W^*(Z),\tau)$
satisfy
\[
 \tau(\xi_i P(Z))=(\tau\otimes\tau)((\partial_i P)(Z)),
\]
where $P$ is a noncommutative polynomial in formal indeterminates
$x_1,\ldots,x_n$, and the difference quotient satisfies
$\partial_i x_j=\delta_{ij}1\otimes1$ and obeys the
Leibniz rule. Put $\Phi^*(Z)=\sum_i\|\xi_i\|_2^2$, or infinity if they
do not exist. We will use the following form of the conjugate-variable
definition.

\begin{lemma}[Polynomial variational formula]\label{lem:variational}
If a bounded self-adjoint tuple $Z$ has conjugate variables, then
\begin{equation}\label{eq:fishervariational}
 \Phi^*(Z)=\sup_{P_1=P_1^*,\ldots,P_n=P_n^*}
 \sum_{i=1}^n\left[
  2(\tau\otimes\tau)((\partial_iP_i)(Z))
  -\tau(P_i(Z)^2)\right],
\end{equation}
where the supremum is over self-adjoint noncommutative polynomials.
It suffices to take $P_i=(Q_i+Q_i^*)/2$ with coefficients of $Q_i$ in
$\mathbb Q+i\mathbb Q$. Consequently, along a family with continuous
moments and existing conjugate variables, free Fisher information is
a measurable function of the parameter.
\end{lemma}

\begin{proof}
The identity
$\partial_i(P^*)=\operatorname{flip}((\partial_iP)^*)$
and the symmetry of $\tau\otimes\tau$ show that $\xi_i^*$ is a
conjugate variable whenever $\xi_i$ is. Uniqueness follows from polynomial
$L^2$ density, so $\xi_i$ is self-adjoint. The expression under the
supremum in \eqref{eq:fishervariational} equals
\[
 \sum_i\big(2\tau(\xi_iP_i(Z))-\|P_i(Z)\|_2^2\big)
 =\sum_i\big(\|\xi_i\|_2^2-\|\xi_i-P_i(Z)\|_2^2\big).
\]
For completeness, the closed span of polynomial vectors is invariant
under each bounded self-adjoint left multiplication operator $Z_j$.
Its orthogonal projection therefore commutes with these operators and
with $W^*(Z)$; since it fixes $1$, it fixes every bounded algebra vector.
Those vectors are dense in $L^2$, proving polynomial density.
Symmetrizing the approximants proves the claimed supremum. Approximating
their finitely many coefficients proves the countable version. Each
functional in that countable supremum is continuous in the moments,
which proves measurability.
\end{proof}

Our nonmicrostates normalization is
\begin{equation}\label{eq:chistar}
 \chi^*(Z)=\frac n2\log(2\pi e)
 +\frac12\int_0^\infty
 \left[\frac n{1+s}-\Phi^*(Z+\sqrt s S)\right]ds.
\end{equation}

Write $\tr_d=d^{-1}\Tr$ and
$\|D\|_{2,d}^2=\sum_i\tr_d(D_i^*D_i)$ for matrix tuples.
A standard GUE tuple $G^{(d)}$ consists of independent centered Gaussian
matrices whose real Hilbert--Schmidt coordinates have covariance $1/d$.
For a random matrix tuple with density $f$, its normalized classical
entropy is
\[
 H_d=-\frac1{d^2}\int f\log f+\frac n2\log d.
\]
Thus independent GUE coordinates of variances $w_i>0$ have normalized
entropy $\frac12\sum_i\log(2\pi e w_i)$.

Voiculescu's asymptotic-freeness theorem introduced the random-matrix
input used here \cite{Vgue}. We use its arbitrary deterministic-tuple
form in \cite[Theorem~22.35]{NS}: if deterministic matrix tuples
converge in normalized trace moments, adjoining independent GUE matrices
gives convergence of expected normalized traces to those of freely
adjoined standard semicirculars. The following lemma supplies the
additional probability and expectation statements needed in this paper.

\begin{lemma}[Bounded matrix inputs]\label{lem:matrix-convergence}
Let $K_d$ be sets of self-adjoint matrix tuples with a common coordinate
operator-norm bound. Suppose their normalized trace moments converge
uniformly over $K_d$ to the moments of a bounded tuple $T$. Let $A^{(d)}$
be any random tuple supported on $K_d$, independent of a fixed finite
standard GUE tuple $G^{(d)}$. Then every normalized polynomial trace in
$(A^{(d)},G^{(d)})$ converges in probability to its value at $(T,S)$,
where $S$ is a standard semicircular family free from $T$. Products of
finitely many such traces converge in expectation as well. The dimensions
may range over any sequence tending to infinity.
\end{lemma}

\begin{proof}
We first record bounds independent of the deterministic input. A
$1/4$-net of the unit sphere of $\C^d$, viewed in real dimension $2d$,
has cardinality at most $9^{2d}$. For a Hermitian matrix $G$,
$\|G\|\leq2\max_v|v^*Gv|$ over this net. For standard GUE, each
$v^*Gv$ is a real centered Gaussian of variance $1/d$, since
$\Tr((vv^*)^2)=1$. Hence
\begin{equation}\label{eq:gue-norm-tail}
 \mathbb P(\|G\|>L)
 \leq 2\exp\big(d(2\log9-L^2/8)\big).
\end{equation}
In particular, the coordinate norms of any fixed finite GUE tuple are
at most $8$ with probability tending to one. For sufficiently large $L$
the right side of \eqref{eq:gue-norm-tail} is at most
$2\exp(-dL^2/16)$. Integrating that tail shows that every fixed moment
of the operator norm is bounded uniformly in $d$.

Now take a deterministic input $a$ with coordinate norms at most $R$,
and put $F=\tr_d P(a,G)$ for a fixed polynomial $P$. Let $(E_\alpha)$
be a Hermitian basis orthonormal for the unnormalized Hilbert--Schmidt
inner product, and write $G_j=\sum_\alpha g_{j\alpha}E_\alpha$.
The real coordinates $g_{j\alpha}$ are independent Gaussians with
variance $1/d$. Cyclicity of trace gives
\[
 \partial_{g_{j\alpha}}F=d^{-1}\Tr(E_\alpha B_j),\qquad
 \sum_{j,\alpha}|\partial_{g_{j\alpha}}F|^2
 =d^{-1}\sum_j\tr_d(B_j^*B_j),
\]
where $B_j$ is the sum of the cyclically ordered products obtained by
deleting each occurrence of the $j$th GUE letter in $P$. The basis also
complexifies to an orthonormal basis of all matrices, so the second
identity applies to complex polynomials.

For a polynomial in standard real Gaussian coordinates, expansion in
orthonormal Hermite polynomials shows that its variance is at most the
expected squared gradient: differentiation weights each nonconstant
coefficient's squared modulus by its positive total degree. Rescaling
the coordinates to variance $1/d$, and applying this to real and
imaginary parts, gives
\begin{equation}\label{eq:gue-trace-variance}
 \E|F-\E F|^2
 \leq d^{-2}\sum_j\E\tr_d(B_j^*B_j)
 =O_{P,R}(d^{-2}).
\end{equation}
The final bound follows from the norm-moment estimate, since each $B_j$
has fixed degree and the deterministic coordinates are uniformly bounded.
Thus the cited expected-trace theorem yields convergence in probability
for each deterministic moment-convergent input sequence.

This convergence is uniform over the sets $K_d$. Otherwise, for some
polynomial and positive error threshold, one could choose a subsequence
and points $a_d\in K_d$ whose error probabilities stay bounded away
from zero. These points form a bounded deterministic moment-convergent
sequence, contradicting the preceding conclusion. Conditioning on
$A^{(d)}$ proves convergence in probability for the random input.

Finally, any fixed product $Q$ of normalized polynomial traces satisfies
\[
 |Q(A^{(d)},G^{(d)})|
 \leq C(1+\max_j\|G_j^{(d)}\|)^D
\]
for constants $C,D$ depending only on the polynomials and the common
cutoff. The same estimate squared has uniformly bounded expectation.
The products are therefore uniformly integrable, so their convergence
in probability implies convergence of their expectations.
\end{proof}
